\section{Introduction}
\label{sec:introduction}

Language models (LMs) are expensive to train. Building a good LM requires a wide range of experimentation over data and architecture choices. Scaling laws, which use the performance of smaller models to predict the performance of a larger model without actually training it, enable more efficient resource allocation for such experimentation.
Further, modern LMs are compared based on their performance on \textit{downstream} tasks, rather than perplexity. 
%\pw{There's some implicit assumption here that existing scaling laws mostly look at perplexity, but this isn't made clear, so it might be confusing to some readers.}
Certain downstream tasks (e.g., HellaSwag) are also often used as development benchmarks for creating training recipes and good data mixtures for training large LMs \citep{OLMo20242O2,Li2024DataCompLMIS}. %\pw{A bit confused by how this sentence is related to the main storyline.}\akshita{making the case for why downstream task prediction is useful}
Scaling laws for downstream tasks are, therefore, important for LM pretraining. 

Prior work has shown results on predicting the average top-1 error over many tasks (as opposed to predicting task accuracy directly) \citep{Gadre2024LanguageMS}, or the accuracy for only one specific task (ARC-Challenge) \citep{Dubey2024TheL3},
using at least 5\% of the compute required to train the target
models. Additionally, these require developing a testbed of many small models to first find compute-optimal models, which adds significant compute costs.
Predicting LM performance on a range of individual downstream tasks in a computationally efficient way (e.g., without finding compute-optimal models) remains an open problem. 

Here, we tackle the challenge of predicting \textbf{individual task performance} of LMs as a function of model size and training data size. In particular, we predict the performance of 7B and 13B models. We focus on a \textbf{range of multiple-choice tasks} and predict the task accuracy for problems written in the ranked classification (RC) format. We use a set of small models (190M to 1.3B non-embedding parameters), which we train for varying durations (1x to 10x Chinchilla optimal data size) to develop scaling laws. Training these \textbf{fixed-size ladder models} costs only 1\% of the compute of the two target models combined. 

We employ a two-step approach to (1) use the number of model parameters $N$ and training tokens $D$ (the \textit{input features}) to predict a task-specific loss (an \textit{intermediate feature}), and then (2) use this task loss to predict \textit{accuracy}. We experiment with different \textbf{design choices} for the intermediate feature, as well as input features, and apply our method on 8 selected tasks from OLMES \citep{gu2024olmes}. We quantify the predictability of the tasks based on the variance of target metrics for a representative small model over the last few training checkpoints. Finally, we provide recommendations for developing downstream scaling laws for new models, and selecting a design choice for a given task. From our recommendations, using a task-specific loss as the intermediate feature, we achieve an absolute error of $<$ 2 points for our target models on four tasks (MMLU, HellaSwag, PiQA, SocialIQA), and an average absolute error of 4 points across both target models and all tasks.

% \ljc{Can we say more tasks? I think we said 4 tasks before. Only saying MMLU doesn't distinguish us from the single-task setting in LLama.} \akshita{yes, I wanted to wait till we rewrite other sections before reporting a final number here}

% Language models (LMs) are expensive to pretrain. Given time and compute constraints, and the variety of decisions that go into building an LM (e.g., picking a ``good'' mixture of pretraining data, making optimal model architecture choices), being able to predict the task performance of a large model before actually training it enables more efficient resource allocation and a wider range of experimentation. This is challenging to do with smaller models ~\citep{Hu2023PredictingEA}.
% Prior work has shown results on predicting the average performance over many tasks \citep{Gadre2024LanguageMS}, or for a single task (ARC-Challenge) \citep{Dubey2024TheL3}, using at least 5\% of the compute required to train the target models.
% Predicting LM performance on individual tasks in a computationally efficient way remains an open problem.

% Here, we tackle the challenge of predicting \textbf{individual task performance} of LMs as a function of model size and training data size. In particular, we predict the performance of 7B and 13B models using 1\% of the compute required to train them. We focus on multiple-choice tasks and predict the task accuracy for problems written in the ranked classification (RC) format. We employ a two-step approach to (1) use the number of model parameters $N$ and training tokens $D$ to predict a task-specific loss (bits-per-byte loss) on the correct answer given the question, then (2) use this task loss to predict accuracy. The closest work \citep{Dubey2024TheL3} only examined compute-optimal models; our approach also works in the \textbf{overtrained regime}, which is important because most recent LMs are heavily overtrained \citep{Dubey2024TheL3,Li2024DataCompLMIS,Bai2023QwenTR}.

% To fit the parameters of the two functions, we measure the task loss and accuracy of a set of small models, varying $N$ between 190M and 1.3B non-embedding parameters, and $D$ between 1x and 10x of the Chinchilla-optimal data size for each model \citep{Hoffmann2022TrainingCL}. Training these \textit{ladder models} costs only 1\% the combined compute of the two target models: a 7B model trained on 4T tokens (7B-4T), and a 13B model trained on 5T tokens (13B-5T). % \footnote{We will release our code and experiments on Github.}

% We apply our method on 8 selected tasks from OLMES \citep{gu2024olmes}.
% On 4 tasks, our predictions are within an absolute error of 2 points for the target models. 
% % Error on the other 4 tasks is higher (average absolute error of 6.9 points) and our analyses show that the metrics of these tasks have relatively high variance between neighboring model checkpoints.
% In \S\ref{sec_variance_analysis}, we quantify the difficulty of the prediction task itself by measuring the variance in the task loss and accuracy of the final few checkpoints (higher variance indicates harder prediction).
% We also show that using less compute for ladder models results in worse predictions (\S\ref{sec:ablate_compute}).
% In \S\ref{sec:ablations}, we ablate various design choices in our prediction method, such as using compute-FLOPs instead of $(N, D)$ as the input, using different intermediate features such as task cross-entropy and general LM loss, and predicting accuracy in a single step. We find that some tasks benefit from these alternatives, but for the tasks that we care the most about (e.g., MMLU and HellaSwag), our two-step approach using task loss works better.
